\section{Limits of the Case-Study Defences}
\label{sec:limits}

StruQ and SecAlign are meaningful improvements over undefended models, but their
reported robustness is conditional on the evaluation design. A low success rate
describes a particular combination of models, tasks, attack construction,
attacker access, optimisation budget, and success rule. It is not a stable
property that can be transferred unchanged to another study or deployment.

This section compares the original case-study evidence with broader and
defence-adaptive evaluations. It also separates security results from utility
measurements and records the limitations of this review.

\subsection{Evaluation Methodology Problems}
\label{subsec:evaluation-methodology}

Attack Success Rate (ASR) is meaningful only when the event counted as success
is explicit. The original StruQ and SecAlign studies primarily test compliance
with a fixed attacker-selected output. StruQ manually inspects responses
containing \enquote{Hacked!} or \enquote{hacked!} to exclude false positives,
whereas SecAlign's default criterion requires the response to begin with
\enquote{Hacked} or \enquote{hacked}. These criteria provide repeatable proxies
for instruction-following compromise, but do not directly measure
confidential-data disclosure, unauthorised tool execution, or the success of
open-ended attacker objectives
\parencites{chen2025struq}{chen2025secalign}{jia2025critical}.

Cross-paper percentages are also limited by prompt diversity and attacker
adaptivity. A defence tested against a fixed set of handcrafted instructions may
not resist an attacker who knows the defence and changes the payload, target, or
optimisation objective. Here, MMLU-PI denotes the study's prompt-injection
benchmark constructed from the Massive Multitask Language Understanding (MMLU)
benchmark
\parencite{jia2025critical}. The empirical studies in
Table~\ref{tab:empirical-comparison} are therefore presented as conditional
evaluations rather than a single ranking.

\begin{landscape}
\scriptsize
\setlength{\tabcolsep}{2pt}
\renewcommand{\arraystretch}{1.15}
\begin{longtable}{@{}p{1.8cm}p{2.2cm}p{1.9cm}p{2.3cm}p{2.1cm}p{1.8cm}p{2.8cm}p{2.5cm}p{1.6cm}p{2.4cm}@{}}
    \caption{Principal empirical evaluations discussed in this paper.}
    \label{tab:empirical-comparison}\\
    \toprule
    \textbf{Defence or evaluation} & \textbf{Model or family} &
    \textbf{Attack type} & \textbf{Attacker access or knowledge} &
    \textbf{Dataset or task} & \textbf{Success metric} &
    \textbf{Reported result} & \textbf{Utility result} &
    \textbf{Source status} & \textbf{Major limitation} \\
    \midrule
    \endfirsthead
    \toprule
    \textbf{Defence or evaluation} & \textbf{Model or family} &
    \textbf{Attack type} & \textbf{Attacker access or knowledge} &
    \textbf{Dataset or task} & \textbf{Success metric} &
    \textbf{Reported result} & \textbf{Utility result} &
    \textbf{Source status} & \textbf{Major limitation} \\
    \midrule
    \endhead

    StruQ original & Llama-7B; Mistral-7B & Manual attacks and adapted GCG &
    GCG uses gradients; trusted task and format known & 208 AlpacaFarm samples
    with a data field & Manually checked ASR for target \enquote{Hacked} &
    Manual attacks \(0\%\)--\(2\%\); defended GCG: Llama \(58\%\), Mistral
    \(56\%\) (undefended \(97\%\), \(99\%\)) & AlpacaEval: Llama
    \(67.2\%\!\to\!67.6\%\); Mistral \(80.0\%\!\to\!78.7\%\) & Peer-reviewed,
    author evaluation
    \parencite{chen2025struq} & Two models; mainly fixed injected task; ASR is a
    proxy for system harm. \\

    SecAlign final & Five Llama/Mistral base and instruct models & Three
    optimisation-free and three optimisation-based attacks & Full instruction
    and format knowledge; white-box optimisation & 208 AlpacaFarm security
    samples; 805 utility prompts & Response begins with \enquote{Hacked} or
    \enquote{hacked} & SecAlign max optimisation-free \(0\%\); optimisation-based
    maxima \(1\%\)--\(14\%\) versus StruQ \(27\%\)--\(60\%\) & AlpacaEval2 win
    rates close to undefended counterparts in the authors' setup & Peer-reviewed,
    author evaluation \parencite{chen2025secalign} & Relative LLM judge, fixed
    default target, and within-study attacker assumptions. \\

    Jia et al. & Released Llama-3-8B-Instruct undefended, StruQ, and SecAlign
    variants & Existing and adaptive Combined/GCG attacks & GCG is white-box;
    adaptive attack knows delimiter defence & OpenPrompt\allowbreak Injection
    and MMLU-PI;
    GCG subsets of 50 and 25 tuples & Task-specific Attack Success Value (ASV) &
    Adaptive GCG: StruQ \(1.00\) on both; SecAlign \(0.46\) on
    OpenPrompt\allowbreak Injection and \(0.88\) on MMLU-PI & Relative win rate: StruQ
    \(21.60\%\), SecAlign \(36.42\%\); absolute
    OpenPrompt\allowbreak Injection/\allowbreak MMLU-PI: StruQ \(0.54/0.35\),
    SecAlign \(0.48/0.34\),
    undefended \(0.65/0.45\) & arXiv v1
    preprint \parencite{jia2025critical} & One principal model family; ASV,
    prompts, and baselines differ from original studies. \\

    Checkpoint-GCG & SecAlign/StruQ Llama, Mistral, Qwen; related
    Meta-SecAlign-8B target & Checkpoint-aware and universal GCG; transfer & Audit
    uses full input and checkpoints; target transfer uses black-box queries &
    AlpacaFarm: 50 randomly selected samples for the exact-input audit; 10
    training and 198 held-out samples for universal evaluation & Target-prefix ASR
    & SecAlign audit: Llama \(88\%\), Mistral \(96\%\), Qwen \(84\%\);
    held-out universal: \(75.3\%\), \(89.9\%\), \(78.3\%\), respectively;
    related-model Llama transfer \(63.9\%\) & Not evaluated & arXiv v2 preprint
    \parencite{yang2025checkpointgcg} & Strong audit access; transfer succeeded
    for one related pair, not the Mistral or Qwen suffixes. \\

    ASTRA and ASTRA++ & SecAlign Llama/Mistral; SecAlign++ & Attention-guided
    white-box optimisation & Weights, gradients, attention, and defence details;
    strong or weak context knowledge & Strong: 40 completed prompts; weak:
    10 train and 198 unseen AlpacaFarm prompts & Response begins with target
    string & At 35 tokens, ASTRA: Llama \(72.5\%\), Mistral \(82.5\%\);
    ASTRA++: SecAlign Llama \(96\%\) (standard deviation 5 percentage points);
    SecAlign++ Llama with the trusted instruction in the system role \(87\%\)
    (standard deviation 12 percentage points) & Not evaluated & arXiv v2
    preprint \parencite{pandya2025astra} & White-box audit, fixed target, larger
    token budget than original evaluations, and limited strong-context sample. \\
    \midrule
    \multicolumn{10}{p{22.5cm}}{\emph{Note:} The rows are not a common ranking.
    Model versions, attack construction, attacker knowledge, optimisation
    budgets, datasets, success rules, and utility baselines differ. ASR and ASV
    should be compared primarily within, rather than across, evaluations.}\\
    \bottomrule
\end{longtable}
\end{landscape}

Utility comparisons require the same care. SecAlign's authors use AlpacaEval2,
where GPT-4-Turbo judges outputs against a \texttt{davinci003} reference. Jia et
al. instead use the corresponding undefended Llama-3-8B-Instruct model as the
relative reference and additionally measure ground-truth task performance on
OpenPromptInjection and MMLU-PI
\parencites{chen2025secalign}{jia2025critical}. The studies therefore answer
different questions and should not be described as directly contradictory.

In the Jia evaluation, StruQ and SecAlign obtain relative win rates of
\(21.60\%\) and \(36.42\%\), respectively, against the undefended instruct model.
Their absolute OpenPromptInjection/MMLU-PI scores are \(0.54/0.35\) for StruQ
and \(0.48/0.34\) for SecAlign, compared with \(0.65/0.45\) for the undefended
model
\parencite{jia2025critical}. These results show utility loss in that setting;
they do not retroactively change the different reference comparison used in the
original SecAlign evaluation.

\subsection{Adaptive Attacks}
\label{subsec:adaptive-attacks}

An adaptive attack targets the assumptions or training process of the defence
being evaluated. It is therefore more informative about worst-case robustness
than a fixed attack suite, although its access requirements may be less
representative of a closed deployment.

Checkpoint-GCG attacks a sequence of intermediate fine-tuning checkpoints. At
each selected checkpoint, GCG is initialised with the suffix found against the
preceding checkpoint; the new suffix is carried forward as a stepping stone
towards the final defended model
\parencite{yang2025checkpointgcg}.

In the exact-input auditing setting, evaluated on 50 randomly selected
AlpacaFarm samples and assuming the full model input and checkpoint access,
Checkpoint-GCG achieved ASRs of \(88\%\), \(96\%\), and \(84\%\) against
SecAlign-defended Llama, Mistral, and Qwen models. Universal
suffixes constructed from ten AlpacaFarm prompts achieved \(75.3\%\), \(89.9\%\),
and \(78.3\%\) on 198 held-out prompts. A Llama-derived suffix also achieved
\(63.9\%\) through black-box queries to the related Meta-SecAlign-8B target
\parencite{yang2025checkpointgcg}. The latter is evidence of transfer for one
closely related model and defence pair; Mistral- and Qwen-derived suffixes did
not transfer in the reported experiment.

ASTRA uses a two-stage white-box search. It first optimises an attention-based
loss to obtain a warm start and then applies a GCG-style target-output objective.
In the strong-context-knowledge evaluation of arXiv v2, the authors completed 40
AlpacaFarm examples using 500 optimisation iterations. At a 35-token budget of
15 prefix and 20 suffix tokens, ASTRA achieved
\(72.5\%\) ASR against SecAlign-defended Llama-3-8B-Instruct and \(82.5\%\)
against SecAlign-defended Mistral-7B-Instruct-v0.1
\parencite{pandya2025astra}.

ASTRA++ withholds the unseen preceding context but retains white-box model and
defence access. Using ten training prompts and 198 held-out prompts, the
35-token configuration of 17 prefix and 18 suffix tokens produced, over three
runs, a mean ASR of \(96\%\) with a standard deviation of 5 percentage points
against SecAlign-defended Llama and \(87\%\) with a standard deviation of 12
percentage points against SecAlign++ when the trusted instruction occupied the
system role
\parencite{pandya2025astra}. These target-prefix results are white-box robustness
audits. They do not establish the same success against a closed API, a dynamic
exfiltration objective, or an arbitrary deployed system.

Checkpoint-GCG and ASTRA demonstrate bypasses of the tested defences under
specified adaptive conditions. Neither study measures a general increase in
attacker effort across deployment settings. The supported conclusion is that
StruQ and SecAlign do not establish a deterministic or independently enforced
boundary in these evaluations, not that every model-level defence must fail.

\subsection{Security and Utility Trade-Offs}
\label{subsec:security-utility}

Defensive training must reduce malicious instruction following without
preventing legitimate analysis of instructional language. A contract-review
task, for example, may need to quote and reason about commands without executing
them. Training a model to distrust every imperative sentence in the data channel
can therefore reduce task utility.

The original StruQ and SecAlign studies report small changes on their relative
LLM-judged benchmarks. The independent Jia evaluation reports lower relative
and absolute utility for released defended Llama variants under its changed
reference and ground-truth tasks. The defensible synthesis is not that one study
invalidates the other, but that utility depends on the model, reference,
benchmark, and task distribution
\parencites{chen2025struq}{chen2025secalign}{jia2025critical}.

Security and utility should consequently be reported together. A defence that
refuses most inputs may record a low attack success rate while being unusable;
a defence that preserves unrestricted instruction following may retain a larger
attack surface. Detector evaluations additionally require false-positive and
false-negative rates rather than aggregate accuracy or area under the curve
(AUC) alone
\parencite{jia2025critical}.

\subsection{Limitations of This Review}
\label{subsec:review-limitations}

This paper is a selective narrative review, not a systematic survey. Its
empirical discussion concentrates on StruQ and SecAlign and on later work that
re-evaluates or adapts attacks to those defences. No original experiments or
independent replications were conducted. Detector- and vendor-specific defences
receive only limited treatment, and closed providers may not disclose the
models, prompts, filters, or updates needed for reproducibility. The adaptive
evidence is predominantly recent preprint work.

Comparison is restricted by changing arXiv versions and by differences in
models, prompts, tasks, attacker budgets, context knowledge, defence knowledge,
and success definitions. Changes to deployed models and configurations also
limit how long an evaluation remains applicable. Target-output ASR is an
incomplete proxy for disclosure or unauthorised tool use. Results are therefore
used to identify conditional patterns, not to construct a universal ranking
\parencite{jia2025critical}.

The review is limited to text-based attacks. It does not evaluate multimodal
injection, persistent memory attacks, denial of service, insiders, or provider
compromise. Stronger architectural controls are also conditional: they depend
on correct trust labels, sufficiently complete policies, trusted interpreters or
security monitors, correct tool wrappers, and an uncompromised trusted computing
base
\parencites{wu2024systemlevel}{debenedetti2025defeating}. The conclusions apply
to the sources and assumptions examined here, not to every possible model-level
defence.

\subsection{Implications of the Evidence}
\label{subsec:open-problem}

The evidence supports a tiered answer to the research question. First, learned
separation substantially reduces the success of the tested handcrafted attacks.
Second, results against standard optimisation-based attacks are mixed but still
show material improvement, particularly for SecAlign in its final evaluation.
Third, later defence-adaptive white-box attacks show that this robustness is not
an independently enforced guarantee. Finally, the transfer of those audit
results to closed, dynamically configured production systems remains unresolved.

Separating instructions and data should therefore be treated as risk reduction,
not complete prevention. Model-level protection remains useful, but deployments
must assume that it can fail and constrain the resulting information flows and
actions outside the model.
